\section{05 Linux Kernel Module}

\subsection{Kernel Modules (LKM)}

%% >>> IMAGE (slide 9): build-up concept -- LKM -> Platform Driver -> + Device Tree -> Access Hardware (GPIO) -> /sys file, driving a 7-segment display <<<

\begin{concept}{Loadable Kernel Module (LKM)}
    A programming module that bridges user/application and hardware. Added
    \important{on the fly} while the OS runs (hence "loadable").
    \begin{itemize}
        \item Not a user program; minimal HW-access functionality
        \item Gateway to user space via a \important{virtual file} (\texttt{/sys/\dots})
        \item Written in \important{C} only (no C++); Rust since Dec 2025
        \item Only \important{root} can load a module
    \end{itemize}
\end{concept}

\begin{example2}{Exam Quiz: Why a kernel driver instead of mmap()? (slide 4)}
    Why implement a driver as a kernel driver if you could access peripheral
    addresses directly via \texttt{mmap()} from user space? Select one or more:
    \begin{enumerate}
        \item A kernel driver can be started automatically when the kernel loads
              the device tree
        \item The GPL license requires using kernel drivers
        \item For self-identifying hardware (USB, PCIe, HDMI) the driver must be in
              the kernel
        \item A kernel driver gives faster access than \texttt{mmap()} from user space
        \item In user space there are no direct interrupts from peripherals
    \end{enumerate}
    \tcblower
    \important{Correct: (1), (3), (5).}
    (2) false -- no GPL requirement; (4) false -- \texttt{mmap()} is \emph{not} slower.
\end{example2}

\mult{2}

\begin{definition}{LKM Commands}
    \begin{itemize}
        \item \texttt{insmod <module.ko>} -- insert
        \item \texttt{rmmod <module>} -- remove
        \item \texttt{lsmod} -- list loaded modules
        \item \texttt{modinfo <module.ko>} -- module info
        \item \texttt{modprobe} -- insert/remove from \texttt{/lib/modules}
    \end{itemize}
\end{definition}

\begin{definition}{Module Types \& Minimum}
    Types: device drivers, filesystem drivers, network drivers. Minimum = an
    \texttt{\_\_init} and an \texttt{\_\_exit} method:
\begin{lstlisting}[language=C, style=basesmol, aboveskip=2pt, belowskip=2pt]
int  __init init_module(void)    { /* init */ }
void __exit cleanup_module(void) { /* clean shutdown */ }
\end{lstlisting}
\end{definition}

\multend

\subsection{Step 1: Simple Kernel Module}

\mult{2}

\begin{definition}{Module Skeleton}
    \begin{itemize}
        \item \texttt{static int \_\_init my\_init(void)}
        \item \texttt{static void \_\_exit my\_exit(void)}
        \item register: \texttt{module\_init(\dots)}, \texttt{module\_exit(\dots)}
        \item info macros: \texttt{MODULE\_LICENSE} / \texttt{AUTHOR} /
              \texttt{DESCRIPTION} / \texttt{VERSION}
    \end{itemize}
\end{definition}

\begin{definition}{Debugging: pr\_* (no printf!)}
    There is \important{no} \texttt{printf}/\texttt{stdio.h} in the kernel.
    \texttt{printk} is deprecated; use:
    \begin{itemize}
        \item \texttt{pr\_info}, \texttt{pr\_err}, \texttt{pr\_alert}, \texttt{pr\_emerg}
        \item \texttt{dev\_info(dev, \dots)}, \texttt{dev\_err(dev, \dots)}
    \end{itemize}
\end{definition}

\multend

\begin{examplecode}{Minimal Kernel Module (slide 13)}
\begin{lstlisting}[language=C, style=basesmol, aboveskip=2pt, belowskip=2pt]
#include <linux/kernel.h>
#include <linux/module.h>
static int __init sevenseg_init(void) {
pr_info("Hello from sevenseg driver!\n");
return 0;
}
static void __exit sevenseg_exit(void) {
pr_info("Bye from sevenseg driver!\n");
}
module_init(sevenseg_init);   // mandatory: identify init
module_exit(sevenseg_exit);   // and cleanup function
MODULE_LICENSE("GPL");
MODULE_AUTHOR("Armin Weiss");
MODULE_DESCRIPTION("7 segment display driver");
MODULE_VERSION("1.0");
\end{lstlisting}
\end{examplecode}

\mult{2}

\begin{corollary}{Tainted Kernel}
    If \texttt{MODULE\_LICENSE} is not a GPL-compatible tag, loading the module
    \important{taints} the kernel.
    \begin{itemize}
        \item \texttt{G}: all modules GPL(-compatible)
        \item \texttt{P}: a proprietary module/kernel is loaded
        \item bit 0 = proprietary, bit 10 = staging, bit 12 = out-of-tree module
        \item check: \texttt{cat /proc/sys/kernel/tainted}
    \end{itemize}
\end{corollary}

\begin{corollary}{Namespace Pollution}
    Without symbol-visibility control, global names in different modules collide
    with the kernel and each other.
    \important{Always use \texttt{static} for global variables and functions!}
\end{corollary}

\multend

\subsection{Step 2: Platform Driver}

\begin{concept}{Discoverable vs. Non-Discoverable HW}
    \begin{itemize}
        \item \important{Discoverable} (PCI, USB): device announces its identity on
              the bus
        \item \important{Non-discoverable} (I2C, SPI): must be declared explicitly in
              the \important{device tree}
        \item A kernel module talking to the device tree = \important{platform driver}
    \end{itemize}
\end{concept}

\begin{definition}{Platform Driver Responsibilities}
    Defines a state struct + a \texttt{probe()} and \texttt{remove()} method.
    \begin{itemize}
        \item \texttt{probe(struct platform\_device *pdev)}: init state, request
              device-tree info, map I/O memory, register IRQs/threads, register to a
              kernel framework
        \item \texttt{remove()}: undo the above
        \item register in \texttt{init()}, unregister in \texttt{exit()}
    \end{itemize}
\end{definition}

\begin{examplecode}{Platform Driver Definition \& Registration (slides 24--25)}
\begin{lstlisting}[language=C, style=basesmol, aboveskip=2pt, belowskip=2pt]
static int sevenseg_probe(struct platform_device *pdev) {
pr_info("Hello from sevenseg_probe()\n"); return 0;
}
static int sevenseg_remove(struct platform_device *pdev) {
pr_info("Bye from sevenseg_remove()\n"); return 0;
}
static struct platform_driver sevenseg_driver = {
.driver = { .name = "sevenseg", .owner = THIS_MODULE },
.probe = sevenseg_probe,
.remove = sevenseg_remove,
};
static int __init sevenseg_init(void) {
return platform_driver_register(&sevenseg_driver);
}
static void __exit sevenseg_exit(void) {
platform_driver_unregister(&sevenseg_driver);
}
\end{lstlisting}
\end{examplecode}

\subsection{Step 3: Add the Device Tree}

\begin{definition}{Matching a Device-Tree Node}
    Provide a \texttt{compatible} list, export it with \texttt{MODULE\_DEVICE\_TABLE},
    and reference it via \texttt{.of\_match\_table}. The driver then binds to the DT
    node with the same \texttt{compatible} string.
\end{definition}

\begin{examplecode}{Device-Tree Match (slide 27)}
\begin{lstlisting}[language=C, style=basesmol, aboveskip=2pt, belowskip=2pt]
static const struct of_device_id sevenseg_ids[] = {
{ .compatible = "sevenseg" },
{ }
};
MODULE_DEVICE_TABLE(of, sevenseg_ids);
static struct platform_driver sevenseg_driver = {
.driver = { .name = "sevenseg", .owner = THIS_MODULE,
.of_match_table = sevenseg_ids },
.probe = sevenseg_probe, .remove = sevenseg_remove,
};
// device tree:  sevenseg { compatible = "sevenseg"; ... };
\end{lstlisting}
\end{examplecode}

\subsection{Step 4: Access the Hardware}

\begin{definition}{Global Data Struct}
    A struct shared between all methods, allocated with \texttt{devm\_kzalloc}
    (auto-freed) and stored via \texttt{platform\_set\_drvdata}.
\begin{lstlisting}[language=C, style=basesmol, aboveskip=2pt, belowskip=2pt]
typedef struct sevenseg_platform_data {
struct device *dev;
void __iomem *base_addr;
int counter, mode, irq;
struct task_struct *kthread_sevenseg;
} sevenseg_platform_data;
// in probe():
pdata = devm_kzalloc(&pdev->dev, sizeof(*pdata), GFP_KERNEL);
platform_set_drvdata(pdev, pdata);
\end{lstlisting}
\end{definition}

\begin{KR}{Access Peripheral Registers from the Driver}
    \textbf{Step 1 -- get the (virtual) base address} from the device tree:
\begin{lstlisting}[language=C, style=basesmol, aboveskip=2pt, belowskip=2pt]
res = platform_get_resource(pdev, IORESOURCE_MEM, 0);  // 35-bit phys
base_addr = devm_ioremap_resource(&pdev->dev, res);    // -> virtual
\end{lstlisting}
    \textbf{Step 2 -- read/write} (8/16/32-bit) at \texttt{base\_addr + offset}:
\begin{lstlisting}[language=C, style=basesmol, aboveskip=2pt, belowskip=2pt]
void iowrite32(u32 value, void *addr);   // write
unsigned int ioread32(void *addr);       // read
\end{lstlisting}
    The DT node carries \texttt{reg = <addr size>} and \texttt{interrupts = <\dots>}.
\end{KR}

\subsection{Step 5: Create a /sys File}

\begin{KR}{Expose a sysfs Attribute}
    \textbf{Step 1 -- show/store callbacks} (read/write of the file):
\begin{lstlisting}[language=C, style=basesmol, aboveskip=2pt, belowskip=2pt]
static ssize_t counter_show(struct device *dev,
struct device_attribute *attr, char *buf) {
return scnprintf(buf, PAGE_SIZE, "value\n");   // user cat
}
static ssize_t counter_store(struct device *dev,
struct device_attribute *attr, const char *buf, size_t count) {
pr_info("New Message %s\n", buf);              // user echo
return count;
}
DEVICE_ATTR(counter, 0644, counter_show, counter_store);
\end{lstlisting}
    \textbf{Step 2 -- create/remove the file} in probe/remove:
\begin{lstlisting}[language=C, style=basesmol, aboveskip=2pt, belowskip=2pt]
device_create_file(&pdev->dev, &dev_attr_counter);   // in probe
device_remove_file(&pdev->dev, &dev_attr_counter);   // in remove
\end{lstlisting}
    \paragraph{conclusion}
    File appears at\\
    \texttt{/sys/bus/platform/drivers/<mod>/<PHYS\_ADDR>.<mod>/counter}\\
    e.g. \texttt{.../sevenseg/fe200000.sevenseg/counter} ($\to$ \texttt{echo}/\texttt{cat}).
\end{KR}

\subsection{Interrupts}

%% >>> IMAGE (slide 44): BCM2711 interrupt routing -- 62 peripheral IRQs through VideoCore + ARM GIC to the 4 ARM cores <<<

\begin{concept}{Interrupt Controller vs. Handler}
    \begin{itemize}
        \item \important{Interrupt controller}: hardware that "executes"/routes the
              interrupt (ARM \important{GIC})
        \item \important{Interrupt handler}: where the ISR is implemented (reacts)
    \end{itemize}
\end{concept}

\begin{definition}{Device-Tree Interrupt Definition}
    \texttt{interrupts = <type number flags>}:
    \begin{itemize}
        \item type: \texttt{0} = SPI (shared), \texttt{1} = PPI (private)
        \item flags: \texttt{1} rising edge, \texttt{2} falling edge,
              \texttt{4} active-high level, \texttt{8} active-low level
        \item controller node: \texttt{compatible = "arm,gic-400"}, referenced via
              \texttt{interrupt-parent}
    \end{itemize}
\end{definition}

\begin{KR}{Calculate the GIC SPI Number (slide 47)}
    For GPIO Bank 0 ($\to$ \texttt{interrupts = <0x00 0x71 0x04>}):
    \begin{enumerate}
        \item \important{VideoCore IRQ}: GPIO Bank 0 (GPIOs 0--27) = IRQ \important{49}
        \item \important{Hardware IRQ}: add ARM offset 96 $\to$ $49 + 96 = 145$
        \item \important{GIC SPI number}: GIC reserves the first 32 (SGI/PPI), so
              subtract 32 $\to$ $145 - 32 = \important{113} = $ \texttt{0x71}
    \end{enumerate}
    So type \texttt{0x00} (SPI), number \texttt{0x71}, flags \texttt{0x04} (active-high level).
\end{KR}

\begin{examplecode}{Register an Interrupt Handler (slide 48)}
\begin{lstlisting}[language=C, style=basesmol, aboveskip=2pt, belowskip=2pt]
static irq_handler_t irq_handler(unsigned int irq, void *dev_id) {
do_something();
iowrite32(0x1 << keys[1],
pdata->base_addr + GPEDSX_OFFSET(keys[1]));  // clear IRQ
return (irq_handler_t) IRQ_HANDLED;
}
// in probe():
pdata->irq = irq_of_parse_and_map(pdata->dev->of_node, 0);
devm_request_irq(pdata->dev, pdata->irq, irq_handler,
IRQF_TRIGGER_NONE, "irq_sevenseg", pdata);
\end{lstlisting}
\end{examplecode}

\subsection{Kernel Threads}

\begin{definition}{kthread API}
    \begin{itemize}
        \item \texttt{kthread\_run(fn, data, namefmt)} -- create + run, returns
              \texttt{task\_struct} (name visible in \texttt{ps})
        \item \texttt{kthread\_should\_stop()} -- check inside the loop
        \item \texttt{kthread\_stop(task)} -- request stop
    \end{itemize}
\end{definition}

\begin{examplecode}{Kernel Thread (slide 52)}
\begin{lstlisting}[language=C, style=basesmol, aboveskip=2pt, belowskip=2pt]
static int thread_fn(void *data) {
while (!kthread_should_stop()) {
do_something();
msleep(2000);
}
return 0;
}
// probe():  pdata->kthread = kthread_run(thread_fn, pdata, "kthread_7seg");
// remove(): kthread_stop(pdata->kthread);
\end{lstlisting}
\end{examplecode}

\subsection{Putting It Together (Lab 5--6)}

\begin{remark}
    \important{7-Segment Counter} combines all three feature types:
    \begin{itemize}
        \item \important{User requests} (write 2-digit numbers) $\to$ sysfs
              (\texttt{counter\_show}/\texttt{counter\_store})
        \item \important{Hardware requests} (start/stop down-counting) $\to$ interrupt
              handler (\texttt{irq\_handler})
        \item \important{Iterative tasks} (periodic down-count) $\to$ kernel thread
              (\texttt{thread\_fn})
    \end{itemize}
    All glued by \texttt{sevenseg\_probe}/\texttt{remove} + the
    \texttt{sevenseg\_platform\_data} struct.
\end{remark}

\begin{remark2}{Memory-Mapped I/O -- Zero-Copy (optional)}
    \texttt{mmap} maps the \important{same physical memory} into both kernel and user
    address spaces $\to$ \important{no copy}, no \texttt{read()}/\texttt{write()}
    syscalls for data transfer.
    \begin{itemize}
        \item driver: \texttt{remap\_pfn\_range()} maps physical pages into the
              user VM; set \texttt{VM\_IO} (excludes from core dumps)
        \item user: \texttt{mmap(\dots, MAP\_SHARED, fd, 0)} $\to$ changes visible to driver
    \end{itemize}
\end{remark2}

\subsection{Exam Exercise (Probepr\"ufung 2021)}

\begin{example2}{Exam 2021 A2: Complete the platform driver}
    Fill in the module so it writes \texttt{1} to the first GPIO address from the
    device tree. Key fill-ins:
    \tcblower
\begin{lstlisting}[language=C, style=basesmol, aboveskip=2pt, belowskip=2pt]
// in probe():
struct resource *res = platform_get_resource(pdev, IORESOURCE_MEM, 0);
void __iomem *base_addr = devm_ioremap_resource(&pdev->dev, res);
iowrite32(1, base_addr);
// match table:
{ .compatible = "brcm,bcm2711-gpio", },
MODULE_DEVICE_TABLE(of, my_module_device_tree_ids);
// driver struct:
.of_match_table = my_module_device_tree_ids,
.probe = my_module_probe, .remove = my_module_remove,
// init / exit:
return platform_driver_register(&my_module_driver);
platform_driver_unregister(&my_module_driver);
module_init(my_module_init); module_exit(my_module_exit);
MODULE_LICENSE("GPL");
\end{lstlisting}
\end{example2}

\raggedcolumns
